\pdfminorversion=7
% =====================================================================
%  Multi-Hop Retrieval-Augmented Generation through the Lens of Evidence Chains
%  Target venue: ACM Transactions on Information Systems (TOIS)
%  Revision v7: strengthened ACM 2026 journal submission format
% =====================================================================

\documentclass[manuscript,review]{acmart}
\citestyle{acmnumeric}

% The acmart review option supplies review line numbering.


% ---------- ACM TOIS submission metadata -----------------------------
\acmJournal{TOIS}
% Volume, issue, article number, month, and DOI are assigned by ACM upon
% acceptance. For initial review, ACM currently requires the single-column manuscript
% format. Journal-specific volume, issue, article number, month, and DOI are
% assigned after acceptance.
\acmDOI{}
\setcopyright{none}
\settopmatter{printacmref=false,printccs=false}
\renewcommand\footnotetextcopyrightpermission[1]{}
% Keep the standard ACM manuscript review page style.

% ---------- Packages -------------------------------------------------
\let\Bbbk\relax
\usepackage{amsmath,amssymb,amsfonts,mathtools}
\usepackage{booktabs}
\usepackage{tabularx}
\usepackage{multirow}
\usepackage{array}
\usepackage{xcolor}
\usepackage{xspace}
\usepackage{etoolbox}
\usepackage{graphicx}
\graphicspath{{figs/}}

% ---------- Figure helpers ------------------------------------------
% The current manuscript uses a corrected evidence-chain figure under ./figs/.
% Figure assets must not contain an internal figure number or caption.
% Use \paperfigure{<width fraction>}{<file>}{<caption>}{<label>}.
\newcommand{\paperfigure}[4]{%
\begin{figure}[tbp]
  \centering
  \includegraphics[width=#1\linewidth]{#2}
  \caption{#3}
  \Description{#3}
  \label{#4}
\end{figure}
}

% Keep floats close to their discussion without creating large white gaps.
\setcounter{topnumber}{2}
\setcounter{bottomnumber}{1}
\setcounter{totalnumber}{3}
\renewcommand{\topfraction}{0.90}
\renewcommand{\bottomfraction}{0.75}
\renewcommand{\textfraction}{0.08}
\renewcommand{\floatpagefraction}{0.80}
\emergencystretch=3em

\usepackage{url}
\usepackage{float}
\usepackage{enumitem}
\usepackage{makecell}
\usepackage{threeparttable}
\usepackage{adjustbox}
\usepackage{rotating}

% ---------- Macros ---------------------------------------------------
\newcommand{\MHRAG}{\textsc{MH-RAG}\xspace}
\newcommand{\RAG}{\textsc{RAG}\xspace}
\newcommand{\LLM}{\textsc{LLM}\xspace}
\newcommand{\PLM}{\textsc{PLM}\xspace}
\newcommand{\KG}{\textsc{KG}\xspace}
\newcommand{\QA}{\textsc{QA}\xspace}
\newcommand{\MHQA}{\textsc{MH-QA}\xspace}
\newcommand{\LLMs}{\textsc{LLM}s\xspace}
\newcommand{\KGs}{\textsc{KG}s\xspace}
\newcommand{\KGQA}{\textsc{KGQA}\xspace}
\newcommand{\LM}{\textsc{LM}\xspace}
\newcommand{\KB}{\textsc{KB}\xspace}
\newcommand{\GNN}{\textsc{GNN}\xspace}
\newcommand{\CoT}{\textsc{CoT}\xspace}
\newcommand{\FiD}{\textsc{FiD}\xspace}
\newcommand{\PRM}{\textsc{PRM}\xspace}
\newcolumntype{Y}{>{\raggedright\arraybackslash}X}
\newcolumntype{C}{>{\centering\arraybackslash}X}

\newcommand{\eunit}{\ensuremath{e}}
\newcommand{\echain}{\ensuremath{\mathcal{C}}}
\newcommand{\kstore}{\ensuremath{\mathcal{K}}}
\newcommand{\retr}{\ensuremath{\mathcal{R}}}
\newcommand{\fuse}{\ensuremath{\mathcal{F}}}
\newcommand{\reason}{\ensuremath{\mathcal{G}}}
\newcommand{\verify}{\ensuremath{\mathcal{V}}}
\newcommand{\pool}{\ensuremath{L_N}}
\newcommand{\budgetset}{\ensuremath{E_K}}
\newcommand{\orderctx}{\ensuremath{\pi(E_K)}}
\newcommand{\obs}{\ensuremath{\mathcal{P}}}
\newcommand{\sel}{\ensuremath{\mathcal{S}}}
\newcommand{\expo}{\ensuremath{\mathcal{O}}}
\newcommand{\fus}{\ensuremath{\mathcal{U}}}
\newcommand{\faith}{\ensuremath{\mathcal{T}}}

\newcommand{\chainset}{\ensuremath{\mathfrak{C}}}

% ---------- Review-scope macro ---------------------------------------
% Update only after the literature and bibliography audit is frozen.
\newcommand{\ReviewCutoff}{July 2026}
\newcommand{\support}{\ensuremath{\operatorname{supp}}}
\newcommand{\doop}{\ensuremath{\operatorname{do}}}
\DeclareMathOperator*{\argmax}{arg\,max}
\DeclareMathOperator*{\argmin}{arg\,min}
\DeclareMathOperator{\rank}{rank}
\DeclareMathOperator{\cost}{cost}
\DeclareMathOperator{\len}{len}
\DeclareMathOperator{\KL}{KL}

\theoremstyle{definition}
\newtheorem{definition}{Definition}
\newtheorem{assumption}{Assumption}
\newtheorem{example}{Example}
\theoremstyle{plain}
\newtheorem{proposition}{Proposition}
\newtheorem{lemma}{Lemma}
\newtheorem{corollary}{Corollary}
\newtheorem{theorem}{Theorem}
\theoremstyle{remark}
\newtheorem{remark}{Remark}
% ---------- Title ----------------------------------------------------
\begin{document}

\title[Multi-Hop RAG through Evidence Chains]
{Multi-Hop Retrieval-Augmented Generation through the Lens of Evidence Chains:\\
A Diagnostic Survey}

% TOIS review is single-blind; author identities remain visible.
\author{Yuqing Luo}
\affiliation{%
  \institution{University of Science and Technology of China}
  \city{Hefei}
  \state{Anhui}
  \country{China}
}
\author{Kai Zhang}
\affiliation{%
  \institution{University of Science and Technology of China}
  \city{Hefei}
  \state{Anhui}
  \country{China}
}

\begin{abstract}
Multi-hop retrieval-augmented generation (RAG) serves information needs whose
support is distributed across evidence units and whose later retrieval steps
depend on intermediate results. This dependency creates an information retrieval
problem beyond ranking several independently relevant documents: a
system must acquire at least one valid evidence chain, preserve it under a
passage, token, or retrieval-call budget, organize it for reader access,
compose its elements correctly, and produce an answer that is faithfully
supported by the supplied evidence. Existing surveys mainly organize the
literature by architecture or control flow, which obscures the stage at which
a method acts and the evidence needed to validate its claim. We present an
IR-centered diagnostic survey organized around five evidence-chain targets:
chain observability, budgeted chain preservation, evidence exposure,
cross-evidence fusion, and causal faithfulness. We synthesize sparse, dense,
iterative, graph-guided, hierarchical, and agentic retrieval; bridge- and
path-aware selection; compression and context organization; fusion and
decomposition methods; and verification and attribution approaches through
\ReviewCutoff. Across these families, we compare the state used to issue the
next query, the unit of selection, the operative budget, the supervision
signal, the reported diagnostic, and the characteristic failure mode. We
then map benchmark annotations and metrics to the targets they can and cannot
identify, distinguish grounded success from ordinary answer accuracy, and
propose a conditional reporting protocol covering retrieval depth, context
budget, membership, ordering, reader capability, faithfulness, and cost. The
framework is intended as a comparative and evaluation tool rather than a
predictive theory of system performance. A versioned companion artifact
records the review protocol, coded entries, and provenance of taxonomy
assignments.
\end{abstract}

\begin{CCSXML}
<ccs2012>
<concept>
<concept_id>10002951.10003317</concept_id>
<concept_desc>Information systems~Information retrieval</concept_desc>
<concept_significance>500</concept_significance>
</concept>
<concept>
<concept_id>10002951.10003317.10003347</concept_id>
<concept_desc>Information systems~Question answering</concept_desc>
<concept_significance>500</concept_significance>
</concept>
<concept>
<concept_id>10010147.10010178.10010179.10003352</concept_id>
<concept_desc>Computing methodologies~Natural language generation</concept_desc>
<concept_significance>300</concept_significance>
</concept>
</ccs2012>
\end{CCSXML}

\ccsdesc[500]{Information systems~Information retrieval}
\ccsdesc[500]{Information systems~Question answering}
\ccsdesc[300]{Computing methodologies~Natural language generation}

\keywords{multi-hop retrieval, retrieval-augmented generation, evidence
chains, evidence selection, context organization, question answering,
faithfulness, evaluation}

\maketitle
\section{Introduction}
\label{sec:introduction}

Retrieval-augmented generation is often described as a two-stage system:
retrieve documents and generate an answer. That description is inadequate
for questions whose evidence is distributed across sources and connected by
dependencies. In a multi-hop question, an answer-bearing passage may not be
retrievable from the original query at all. It becomes identifiable only
after an earlier passage reveals a bridge entity, relation, date, or partial
conclusion. The retrieval problem is therefore sequential and conditional:
the system must discover a support structure, not merely collect several
high-scoring documents.

This distinction matters for both system design and evaluation. A complete
support chain may be absent from the initial retrieval pool. It may be
present in the pool but removed by a top-$K$ selector that favors marginal
query relevance. It may survive selection but be placed in a reader-unfriendly
position, separated from its companion evidence by distractors, or compressed
in a way that removes the bridge relation. Even when all required evidence is
available, the reader may fail to combine it. Conversely, a correct answer
may be produced from parametric memory, annotation artifacts, or a single-hop
shortcut without using the retrieved chain. End-to-end exact match or F1
collapses these cases into one number and therefore gives limited guidance
about where a method succeeds.

For a question $q$, let $\chainset^\star(q)$ denote the set of valid minimal
evidence chains that support the reference answer. A retrieval system
constructs a candidate pool $L_N$, a selector places a subset $E_K$ under a
passage or token budget, an ordering or packing policy produces
$\pi(E_K)$, and a reader generates $\hat a$. We use this sequence to organize
the literature around five diagnostic questions:
\begin{enumerate}[leftmargin=1.5em]
\item Does $L_N$ contain at least one valid chain?
\item Conditional on chain presence, does the selected context preserve one?
\item Conditional on preservation, is the chain presented where and how the
reader can access it?
\item Given accessible evidence, can the reader perform the required
cross-evidence composition?
\item Is the resulting answer causally supported by the supplied evidence?
\end{enumerate}
These questions correspond to chain observability, selection preservation,
evidence exposure, fusion reliability, and causal faithfulness. They are not
independent modules, and a single system change can affect several of them.
Their value is diagnostic: they separate claims that are routinely conflated
in multi-hop RAG experiments.

\paperfigure{1.0}{fig1_evidence_chain_v3.pdf}
{An evidence-chain view of multi-hop RAG. A system constructs a retrieval
pool $L_N$, selects a budgeted context $E_K$, presents the context through an
ordering or packing policy $\pi$, and generates an answer. Evaluation should
separately test whether a valid chain is observable, preserved, exposed,
correctly fused, and causally used. Ordinary answer accuracy also includes
answers obtained through shortcuts, parametric memory, or guessing.}
{fig:latent_chain_view}

The product shown in Figure~\ref{fig:latent_chain_view} is the exact
chain-rule decomposition of the grounded-success event $Y_g$ under the declared
evaluation protocol; it is not an approximation to ordinary answer accuracy.

Existing surveys provide broad taxonomies of RAG architectures, multi-hop
question answering, graph reasoning, chain-of-thought, and agentic
retrieval~\cite{gao2023ragsurvey,fan2024ragllmsurvey,
mavi2022mhqasurvey,mohammadi2024mhmrcsurvey,
plaat2024multistepsurvey,singh2025agenticragsurvey,
liang2025reasoningragsurvey}. These taxonomies answer an important procedural
question: how is a system built? They are less suited to a second question:
which failure event does the central mechanism change? An iterative dense
retriever, a graph expander, and a search agent have different
architectures, yet all can increase the probability that a bridge-dependent
chain enters the candidate pool. A cross-encoder, a bridge-aware selector,
and a compressor can all change which chain elements survive the final
budget. A long-context model and an ordering policy can both affect evidence
access without altering retrieval recall.

The survey therefore uses two complementary views. The primary view follows
the evidence chain from acquisition to causal use. A compact architectural
crosswalk is retained to connect this organization to familiar method
families. This positioning is deliberately narrower than a general RAG or
LLM-reasoning survey. We focus on systems in which external evidence and
cross-evidence dependency are necessary parts of the task. Purely parametric
reasoning, general-purpose agents without retrieval-grounded evidence use,
and single-hop retrieval methods are discussed only when they provide a
component or baseline for multi-hop systems.

\subsection{Position Relative to Existing Surveys}
\label{subsec:survey_position}

Prior work surveys general RAG~\cite{gao2023ragsurvey,fan2024ragllmsurvey,
huang2024ragsurvey,zhao2024ragaigcsurvey,li2024ragdecoupledsurvey}, multi-hop
question answering and reading comprehension~\cite{mavi2022mhqasurvey,
mohammadi2024mhmrcsurvey}, complex and knowledge-graph QA~\cite{lan2021complexkbqasurvey,
gu2022knowledgesurvey,pan2024llmkgsurvey}, step-by-step reasoning~\cite{plaat2024multistepsurvey,
chu2024cotsurvey,qiao2023reasoningsurvey}, attribution and hallucination~\cite{li2023attributionsurvey,
huang2023hallucinationsurvey}, and agentic or reasoning-oriented RAG~\cite{singh2025agenticragsurvey,
liang2025reasoningragsurvey}. A close process-centric framework organizes
retrieval--reasoning systems by execution plan, index structure, next-step
control, and stopping criterion~\cite{ji2026retrievalreasoning}. These views
explain what systems do. Our complementary view asks which evidence-chain
failure stage a mechanism changes and which conditional evidence supports
that claim.

\begin{table}[t]
\caption{Positioning of the present survey relative to major survey families.}
\label{tab:survey_crosswalk}
\centering
\scriptsize
\begin{tabularx}{\linewidth}{@{}p{2.5cm}p{2.7cm}p{2.8cm}Y@{}}
\toprule
Survey family & Representative scope & Primary organizing axis & Gap addressed here \\
\midrule
General RAG~\cite{gao2023ragsurvey,fan2024ragllmsurvey} & Retrieval, augmentation, and generation & System architecture and lifecycle & Multi-hop failure stages are usually aggregated. \\
Multi-hop QA / MRC~\cite{mavi2022mhqasurvey,mohammadi2024mhmrcsurvey} & Bridge, comparison, and compositional reasoning & Task and reasoning type & Pool, budget, presentation, and evidence use are not separated. \\
Reasoning and agentic RAG~\cite{singh2025agenticragsurvey,liang2025reasoningragsurvey} & Adaptive retrieval, planning, and tool use & Control flow and agent design & End-to-end gains remain difficult to attribute. \\
Retrieval--reasoning process~\cite{ji2026retrievalreasoning} & Plan, index, control, and stopping & Procedural decisions & No stage-conditional metric and reporting map. \\
\textbf{This survey} & \textbf{Multi-hop RAG over text, graphs, tables, and hybrid sources} & \textbf{Evidence-chain stages plus architecture} & \textbf{Conditional diagnostics, failure attribution, and reporting protocol.} \\
\bottomrule
\end{tabularx}
\end{table}

The ingredients of this framing are not new in isolation. Retrieved documents
have been modeled as latent variables in end-to-end RAG~\cite{lewis2020rag,guu2020realm};
multi-hop datasets annotate support chains~\cite{yang2018hotpotqa,ho2020wikimultihopqa,trivedi2022musique};
and attribution work studies whether answers depend on evidence~\cite{gao2023alce,liu2023evaluating,wu2024cofca,xu2025veritas}.
The contribution is their integration into one IR-centered diagnostic
organization that supports cross-architecture comparison and exposes what
common metrics do not identify.

The survey is guided by four research questions:
\begin{description}[leftmargin=1.5em,style=nextline]
\item[RQ1: Evidence acquisition.] Which retrieval mechanisms make dependent
support chains observable, and how do they trade coverage against noise,
cost, and query drift?
\item[RQ2: Budget and presentation.] How do selection, compression, ordering,
and context construction determine whether an observable chain remains
usable by the reader?
\item[RQ3: Composition and support.] Which reader, decomposition, graph, and
verification mechanisms improve cross-evidence reasoning and evidence-grounded
answers?
\item[RQ4: Evaluation.] Which benchmarks, annotations, and interventions can
identify each failure stage, and which common metrics conflate them?
\end{description}

\subsection{Contributions}
\label{subsec:contributions}

This survey makes four contributions.
\begin{itemize}[leftmargin=1.3em]
\item \textbf{An IR-centered evidence-chain framework.} We formalize the
observable sequence from retrieval pool to generated answer and define five
protocol-bound events that make retrieval, selection, presentation,
reasoning, and faithfulness claims comparable across architectures.
\item \textbf{A dual-axis review.} We organize methods primarily by the stage
they target while retaining an architectural crosswalk covering iterative
retrieval, graph-based methods, decomposition, fusion readers, long-context
systems, verification, and agentic search.
\item \textbf{Benchmark and metric alignment.} We reinterpret supporting-fact,
path-annotated, long-context, heterogeneous-evidence, and counterfactual
benchmarks according to the failure events they expose, and we distinguish
pool recall, context preservation, fixed-membership presentation, oracle-
evidence reasoning, and causal-support diagnostics.
\item \textbf{A reporting protocol and auditable artifact.} We propose a
minimum set of conditional diagnostics and budget measures for multi-hop RAG
studies, and release the literature protocol and coded records with explicit
verification status.
\end{itemize}

\paragraph{Organization.}
Section~\ref{sec:survey_methodology} describes the review protocol and scope.
Sections~\ref{sec:problem_setup} and~\ref{sec:stat_decomposition} introduce
the task model and diagnostic decomposition. Section~\ref{sec:estimand_methods}
reviews method families along the evidence-chain pipeline, and
Section~\ref{sec:architecture_view} provides the architectural crosswalk.
Section~\ref{sec:benchmarks} aligns benchmarks and metrics with diagnostic
claims. Section~\ref{sec:causal_faithfulness} focuses on causal evidence use.
Sections~\ref{sec:tradeoffs} and~\ref{sec:reporting_guidelines} summarize open
problems and reporting recommendations.
\section{Review Protocol and Literature Coverage}
\label{sec:survey_methodology}

The review is designed as a structured evidence map rather than a statistical
meta-analysis. Its purpose is to compare mechanisms and evaluation claims
across communities whose terminology differs, while preserving a traceable
record of why each work is included and how it is coded.

\subsection{Search Sources and Time Window}
\label{subsec:search_protocol}

We searched ACL Anthology, OpenReview, arXiv, ACM and other publisher
proceedings pages, and official project repositories. DBLP and Semantic
Scholar were used for secondary discovery and deduplication. The review
covers work from January 2017 through July 2026. Query families combined task
terms (e.g., \emph{multi-hop question answering}, \emph{multi-step
retrieval}, \emph{compositional question answering}), mechanism terms
(e.g., \emph{iterative retrieval}, \emph{graph retrieval}, \emph{question
decomposition}, \emph{context compression}, \emph{agentic search}), and
evaluation terms (e.g., \emph{supporting facts}, \emph{path recall},
\emph{citation faithfulness}, \emph{counterfactual evaluation}). Backward
and forward citation tracing from major surveys, benchmarks, and canonical
methods supplemented keyword search.

\subsection{Inclusion, Exclusion, and Coding}

A work is included when it satisfies at least one of four conditions: it
proposes a method for retrieving or using dependent evidence; introduces a
benchmark or annotation scheme that exposes multi-step evidence dependency;
analyzes evidence use or failure in a retrieval-grounded setting; or provides
a foundational component repeatedly used by multi-hop RAG systems. We exclude
works whose multi-hop framing is incidental, general agents without a
retrieval-grounded evidence objective, and single-hop methods that do not
materially inform a multi-hop stage.

Each method record stores its evidence substrate, retrieval setting, context
budget, architectural family, primary diagnostic target, secondary targets,
benchmarks, cost information, and source-verification status. The primary
target is the earliest evidence-chain stage directly changed by the central
mechanism; later effects are recorded as secondary targets. For example, a
query-rewriting method that expands the candidate pool is coded primarily as
observability, even when the resulting evidence also improves fusion.

The coding scheme separates three objects that should not be conflated:
\emph{discovery records}, which maximize recall during literature search;
\emph{reviewed method and benchmark records}, whose metadata and assignments
have been checked against the paper or official project page; and
\emph{manuscript citations}, which additionally include background and
component papers. Only reviewed records support taxonomy-level comparisons.
The companion artifact reports these layers and their status independently.

\subsection{Quality Control and Limitations}

Borderline assignments are logged with the competing labels, the relevant
source passage, and the reason for the final decision. Records whose primary
intervention stage remains unresolved after full-text checking stay in a
\emph{needs-review} state and are excluded from taxonomy-level comparisons.
Because most records are single-coded rather than independently double-coded,
we do not report an inter-annotator agreement coefficient. Absence from the
catalog is likewise not interpreted as evidence of irrelevance or low
quality.

Three limitations follow. First, rapid publication in 2025--2026 means that
titles, venue status, and versions can change; the artifact therefore records
a cut-off date and canonical source. Second, valid support chains are often
non-unique, so benchmark annotations identify only a subset of the chains a
system might use. Third, bibliography breadth is not equivalent to verified
taxonomy coverage. The companion artifact therefore records a versioned
release, the reviewed-record denominator used by generated comparisons, and
the common source-level checklist required for every included entry.

\subsection{Quantitative Reporting Audit}
\label{subsec:quantitative_audit}

To avoid treating bibliography size as empirical coverage, we conducted a
small reporting audit on the frozen reviewed seed used for this revision. The
denominator is the set of method records in the companion artifact with
source-verified status, excluding discovery-only and needs-review entries.
The audit is therefore a diagnostic sample of the catalog, not a frequency
estimate over all multi-hop RAG publications.

\begin{table}[t]
\caption{Frozen seed audit over source-verified method records. Counts use
the reviewed denominator $n=20$ and do not include discovery-only entries.}
\label{tab:reviewed_seed_audit}
\centering
\scriptsize
\begin{tabularx}{\linewidth}{@{}p{3.1cm}p{2.0cm}Y@{}}
\toprule
Audit item & Count & Interpretation \\
\midrule
Primary target: observability & 7/20 & Acquisition remains the dominant
coded intervention target. \\
Primary target: budgeted selection & 3/20 & Selection-preservation claims
are present but less often isolated than acquisition claims. \\
Primary target: evidence exposure & 2/20 & Context organization appears
mainly in hierarchical or graph-summary systems. \\
Primary target: fusion reliability & 5/20 & Reader and graph-reasoning
methods remain a substantial family. \\
Primary target: causal faithfulness & 0/20 & No reviewed seed method is
primarily coded as a causal-use intervention. \\
Joint or adaptive primary control & 3/20 & Agentic systems often span
several events and require event-specific logs before attribution. \\
Public code URL recorded & 13/20 & Reproducibility metadata is common but
not universal in the reviewed seed. \\
Benchmark records reviewed & 6 & Most benchmark entries identify joint
answer or support behavior rather than a single controlled stage. \\
\bottomrule
\end{tabularx}
\end{table}

Two patterns are visible even in this conservative seed. First, the catalog is
front-loaded toward acquisition and fusion: observability and fusion account
for 12 of the 20 reviewed method records, whereas selection and exposure are
usually evaluated through downstream answer gains rather than conditional
chain-preservation or fixed-membership tests. Second, faithfulness is mostly
a secondary or evaluation-side concern in the seed. This supports the
reporting protocol in Section~\ref{sec:reporting_guidelines}: future catalog
releases should code whether each paper reports stage-aligned metrics, token
budget, retrieval-call budget, latency, membership/order controls, and
deletion, conflict, or counterfactual evidence-use probes. Those fields are
required before stronger temporal or architecture-family reporting-gap claims
can be made.
\section{Problem Setup: Multi-Hop Retrieval as Evidence-Chain Acquisition}
\label{sec:problem_setup}

\subsection{What Makes Retrieval Multi-Hop?}
\label{subsec:what_is_mh}

A retrieval task is multi-hop when at least one required evidence unit becomes
retrievable or interpretable only after an intermediate state has been
resolved. This definition emphasizes dependency rather than the raw number of
documents. A question can require many passages but remain single-hop if all
passages are independently query-relevant. Conversely, a two-document bridge
question is multi-hop when the second document cannot be specified before the
first reveals an entity or relation.

Three properties are useful in practice. \emph{Semantic non-locality} means
that some required evidence is weakly related to the original query surface.
\emph{Dependency} means that an intermediate result changes the next
retrieval or interpretation step. \emph{Joint sufficiency} means that no
proper subset of the selected support is sufficient under the task protocol.
These properties distinguish multi-hop RAG from long-document reading,
ordinary multi-document aggregation, and difficult single-evidence reasoning.

\subsection{Evidence Chains and Observable System Objects}
\label{subsec:latent_chain}

Let $q$ be a question, $a^\star$ its reference answer, and $\kstore$ an
external knowledge store. An evidence unit may be a sentence, passage, table
row, image region, knowledge-graph triple, or another retrievable object. Let
$\chainset^\star(q)$ denote the set of valid minimal evidence chains for $q$.
A chain is sufficient if it supports $a^\star$ under the task definition and
minimal if removing any required unit destroys that sufficiency. The set-valued
definition accommodates alternative valid sources and paths.

A system exposes the following observable sequence:
\begin{align}
L_N &= \retr_N(q,\kstore), \\
E_K &= S_K(q,L_N), \qquad E_K\subseteq L_N, \\
\pi(E_K) &= O(q,E_K), \\
\hat a &= \reason(q,\pi(E_K)),
\end{align}
where $L_N$ is the retrieval pool, $E_K$ is a selected context under token or
passage budget $K$, $\pi$ is the reader-facing ordering or packing policy,
and $\hat a$ is the generated answer. We use $K$ as a generic context budget rather than only a passage
count because compression and heterogeneous evidence make unit lengths
variable.

\subsection{Running Example}
\label{subsec:running_example}

Consider the HotpotQA question: \emph{``What government position was held by
the woman who portrayed Corliss Archer in the film Kiss and Tell?''}
The first support unit identifies Shirley Temple as the actor who portrayed
Corliss Archer; the second states that she later served as United States Chief
of Protocol~\cite{yang2018hotpotqa}. The second unit is poorly specified by
the original query and becomes straightforward to retrieve only after the
bridge entity has been resolved.

This example separates five failures. The biography may be missing from
$L_N$; it may appear in the pool but be removed from $E_K$; both passages may
be selected but placed or compressed so that the reader cannot connect them;
the reader may fail to compose the cast and biographical facts; or it may
produce the correct title from parametric memory without using the supplied
chain. The same answer-level error can therefore originate at different
stages, and a correct answer does not certify retrieval-grounded reasoning.

\subsection{Adjacent Problems and Scope Boundaries}

Graph paths, question decompositions, and agent trajectories are possible
representations of an evidence chain, but none is identical to the chain by
definition. A graph system can retrieve a path that is not sufficient; an
agent can execute several tool calls without acquiring dependent evidence; a
chain-of-thought can be multi-step while remaining entirely parametric. This
survey includes these systems only when their external evidence acquisition,
selection, composition, or verification is central to the task.

Long-context systems also require care. Enlarging the input may reduce the
need for repeated retrieval, but it does not remove dependency, budget, or
position effects. Similarly, multimodal and table-text systems change the
type of evidence unit while preserving the same acquisition and composition
problem. These substrates are discussed where they reveal a distinct
retrieval or fusion mechanism rather than as independent application domains.
\section{A Diagnostic Decomposition of Grounded Success}
\label{sec:stat_decomposition}

The evidence-chain view is used as an evaluation decomposition. It does not
assume that components are independent, that the stages can be changed one at
a time, or that a method has an intrinsic score independent of its corpus,
budget, reader, and intervention protocol.

\subsection{Five Stage Events}

For a question $q$, define the following events:
\begin{align}
\obs &= \{\exists C\in\chainset^\star(q): C\subseteq L_N\}, \\
\sel &= \{\exists C\in\chainset^\star(q): C\subseteq E_K\}, \\
\expo &= \{\exists C\in\chainset^\star(q): C\subseteq E_K,
          \operatorname{access}_{\pi}(C)\le r\}, \\
\fus &= \{\hat a=a^\star \text{ under the realized chain-containing,
          accessibility-controlled context}\}, \\
\faith &= \{\exists C\in\chainset^\star(q): C\subseteq E_K,
          \hat a=a^\star,\; \hat a\leftarrow C\}.
\end{align}
The accessibility function is protocol-specific. It can encode the latest
support position, the maximum separation between dependent units, or the
outcome of a controlled packing treatment. It must therefore be declared for
the reader and context format being evaluated. The notation
$\hat a\leftarrow C$ denotes the latent claim that the answer causally depends
on a valid chain $C$; it is not directly observed from answer text or citation
overlap.

The grounded-success event is
\begin{equation}
Y_g=\obs\cap\sel\cap\expo\cap\fus\cap\faith.
\end{equation}
Applying the chain rule in pipeline order gives
\begin{align}
P(Y_g)=&\;P(\obs)P(\sel\mid\obs)
P(\expo\mid\obs,\sel) \\
&\times P(\fus\mid\obs,\sel,\expo)
P(\faith\mid\obs,\sel,\expo,\fus).
\label{eq:success_decomp}
\end{align}
Equation~\ref{eq:success_decomp} is an exact bookkeeping identity once the
population and evaluation protocol are fixed. It is not a product model for
predicting end-to-end performance. Its practical role is to specify the
conditional subset on which each claim should be tested. Ordinary answer
success $Y=\{\hat a=a^\star\}$ also contains $Y\cap\neg Y_g$, including
parametric-memory answers, single-hop shortcuts, incomplete but lucky
reasoning, and guessing.

\begin{table}[t]
\caption{Evidence-chain stages, suitable diagnostics, and common confounds.}
\label{tab:diagnostic_protocol}
\centering
\small
\begin{tabularx}{\linewidth}{@{}p{2.4cm}p{3.1cm}p{3.1cm}Y@{}}
\toprule
Stage & Target quantity & Suitable diagnostic & Main confound \\
\midrule
Pool construction & $P(\obs)$ & complete-chain or path recall at Top-$N$ & incomplete or non-unique gold chains \\
Budgeted selection & $P(\sel\mid\obs)$ & chain preservation under a fixed token or passage budget & changing retrieval depth or evidence granularity \\
Presentation & $P(\expo\mid\obs,\sel)$ & fixed-membership ordering or packing ablation & changing membership and order together \\
Fusion & $P(\fus\mid\obs,\sel,\expo)$ & oracle-evidence or controlled-context EM/F1; step validity & parametric answer availability \\
Causal faithfulness & $P(\faith\mid\obs,\sel,\expo,\fus)$ & deletion, conflict, and bridge-counterfactual tests & redundant chains and out-of-distribution interventions \\
\bottomrule
\end{tabularx}
\end{table}

\subsection{Conditional Evidence Contribution}
\label{subsec:conditional_utility}

Query--document relevance is an incomplete selection signal in a dependent
retrieval task. For a fixed reader and partial context $C$, a useful mechanism
quantity is the conditional contribution
\begin{equation}
u(e\mid q,C)=F_q(C\cup\{e\})-F_q(C),
\end{equation}
where $F_q(C)$ is the reader's probability of the correct answer under
context $C$. This signed, reader-dependent quantity is not a stage probability
and is generally expensive to estimate. It nevertheless clarifies why a
bridge passage can have low query-only relevance and high value after a core
passage is known. Bridge-aware reranking, path scoring, note construction,
and leave-one-out supervision estimate different surrogates of $u$.

\subsection{Budget, Exposure, and Fusion}
\label{subsec:fusion_program}

The transition from $L_N$ to $E_K$ is a constrained set-selection problem.
A deeper pool can increase observability while adding distractors; a smaller
reader context can reduce noise while discarding a complementary bridge.
Compression changes both the cost and the semantic content of an evidence
unit and should therefore be evaluated for chain preservation and provenance,
not only token reduction.

Presentation must be separated from membership. If a reranker changes both
which passages are selected and where they appear, an answer gain cannot be
attributed to exposure. The cleanest test fixes $E_K$ and varies only
$\pi(E_K)$. Fusion should likewise be evaluated on examples where a valid
chain is supplied in a controlled context. Oracle-evidence accuracy is not a
deployment score; it is a reader diagnostic for bridge resolution,
comparison, aggregation, or verification once evidence acquisition has been
removed as a failure source.

\subsection{Identifiability under Partial Annotation}

Supporting-fact and path annotations identify only annotated versions of the
pool, selection, and exposure events. A positive match establishes that at
least one labeled chain is present; a negative match does not rule out an
unannotated valid chain. Fusion and causal faithfulness are less identifiable
from observational outputs. A correct answer with a plausible citation can
still be produced from parametric memory. Evidence removal, factual conflict,
and bridge substitution provide stronger evidence, but deletion can create
unnatural inputs and alternative valid chains can mask sensitivity. Stage
metrics should therefore be reported as protocol-bound estimators with their
annotation and intervention assumptions made explicit.

\section{Method Families Along the Evidence-Chain Pipeline}
\label{sec:estimand_methods}

We code a method by the earliest evidence-chain stage directly changed by its
central mechanism and record later effects as spillovers. This rule prevents
an architectural label from standing in for a causal claim: graph methods can
alter acquisition, selection, or fusion; decomposition can generate retrieval
queries or simplify reader computation; and an agent can act across the full
pipeline. Each subsection therefore asks both what state the method uses and
what evidence would be required to validate its stage-aligned claim.

\subsection{Evidence Acquisition and Chain Observability}
\label{subsec:methods_observability}

Acquisition methods differ mainly in the state used to issue the next query,
the unit expanded, and the rule that stops search. Table~\ref{tab:acquisition_design}
summarizes these distinctions.

\begin{table}[t]
\caption{Design choices for evidence acquisition. ``Diagnostic'' denotes the
minimum evidence needed to support an observability claim.}
\label{tab:acquisition_design}
\centering
\scriptsize
\begin{tabularx}{\linewidth}{@{}p{2.25cm}p{2.6cm}p{2.3cm}p{2.5cm}Y@{}}
\toprule
Family & State for the next query & Search unit / budget & Diagnostic & Characteristic risk \\
\midrule
Single-shot sparse or dense & original question & passages; fixed Top-$N$ & complete-chain recall at fixed $N$ & later hops absent from the query surface \\
Iterative / interleaved & question plus retrieved or generated state & passages; hop or call budget & chain recall by hop and call count & query drift and error propagation \\
Graph expansion / traversal & local entities, paths, or subgraph state & nodes, edges, linked text; expansion depth & path or chain recall plus graph coverage & missing links, noisy linking, explosive neighborhoods \\
Hierarchical / coarse retrieval & summaries, clusters, or hypothetical units & variable-granularity units; token budget & chain coverage at matched reader tokens & compression and granularity confounds \\
Agentic / RL search & reasoning history, tool observations, memory & queries and actions; learned stopping & chain recall, actions, tokens, and latency & more compute and weak action attribution \\
\bottomrule
\end{tabularx}
\end{table}

\paragraph{Single-shot pool construction.}
BM25~\cite{robertson2009bm25}, learned sparse retrieval such as
\textsc{SPLADE}~\cite{formal2021splade}, dense retrievers including
\textsc{DPR}~\cite{karpukhin2020dpr}, \textsc{REALM}~\cite{guu2020realm},
\textsc{ANCE}~\cite{xiong2020ance}, \textsc{Contriever}~\cite{izacard2022contriever},
\textsc{GTR}~\cite{ni2022gtr}, \textsc{E5}~\cite{wang2022e5},
\textsc{BGE}~\cite{xiao2024bge}, and \textsc{RetroMAE}~\cite{xiao2022retromae},
and late-interaction systems such as \textsc{ColBERT} and
\textsc{ColBERTv2}~\cite{khattab2020colbert,santhanam2022colbertv2}
provide scalable candidate pools. They are indispensable baselines, but their
query state is fixed. When a later hop is specified only by a bridge revealed
inside earlier evidence, changing the matching function alone does not remove
the structural information gap.

\paragraph{Iterative and reasoning-interleaved retrieval.}
\textsc{GoldEn Retriever}~\cite{min2019golden}, \textsc{MDR}~\cite{xiong2021mdr},
\textsc{Baleen}~\cite{khattab2021baleen}, and \textsc{MUPPET}~\cite{feldman2019muppet}
condition later retrieval on accumulated evidence or bridge state.
\textsc{IRCoT}~\cite{trivedi2023ircot} uses intermediate reasoning as a
query state; \textsc{FLARE}~\cite{jiang2023flare},
\textsc{ITER-RETGEN}~\cite{shao2023iterretgen},
\textsc{Adaptive-RAG}~\cite{jeong2024adaptiverag},
\textsc{DRAGIN}~\cite{su2024dragin}, \textsc{Self-RAG}~\cite{asai2024selfrag},
and \textsc{ITER-DRAGON}~\cite{wu2024iterdragon} adapt retrieval timing,
depth, or filtering to generator state. Their common benefit is state-
conditioned reachability; their common confound is that more calls also
change compute, pool size, and downstream context.

\paragraph{Graph and hierarchical acquisition.}
Question-specific graph systems such as \textsc{GRAFT-Net}~\cite{sun2018graftnet},
\textsc{PullNet}~\cite{sun2019pullnet}, \textsc{HGN}~\cite{fang2020hgn}, and
\textsc{DFGN}~\cite{qiu2019dfgn} expand through entity--passage structure.
\textsc{ToG}, \textsc{ToG 2.0}, and \textsc{RoG}~\cite{sun2024tog,ma2024tog2,luo2024rog}
traverse knowledge-graph paths under language-model control, while
\textsc{KG-FiD} and \textsc{GNN-RAG}~\cite{yu2022kgfid,mavromatis2024gnnrag}
use graph signals to prune or aggregate retrieved evidence. Corpus-level graph
indexes such as \textsc{GraphRAG}, \textsc{LightRAG}, and
\textsc{HippoRAG}~\cite{edge2024graphrag,guo2024lightrag,gutierrez2024hipporag}
change the retrieval substrate itself. Hierarchical methods including
\textsc{RAPTOR}, \textsc{LongRAG}, and \textsc{HiQA}~\cite{sarthi2024raptor,jiang2024longrag,chen2024hiqa}
change evidence granularity, while \textsc{HyDE}~\cite{gao2023hyde} changes
the query representation through a hypothetical document. For these systems,
path reachability or summary coverage should be reported separately from the
reader gain produced by a coarser context.

\paragraph{Agentic and adaptive search.}
\textsc{ReAct} and \textsc{Toolformer}~\cite{yao2023react,schick2023toolformer}
made search and tool invocation part of the reasoning trajectory. Recent
systems span reinforcement-learned search agents---including
\textsc{Search-R1}~\cite{jin2025searchr1},
\textsc{R1-Searcher}~\cite{song2025r1searcher},
\textsc{ReSearch}~\cite{chen2025research}, and
\textsc{RAG-Gym}~\cite{xiong2025raggym}---and adaptive or structured
controllers such as \textsc{DeepRAG}~\cite{guan2025deeprag},
\textsc{PAR-RAG}~\cite{zhang2025parrag},
\textsc{DualRAG}~\cite{cheng2025dualrag},
\textsc{Search-o1}~\cite{li2025searcho1}, and
\textsc{MultiCube-RAG}~\cite{shi2026multicuberag}. These systems place query
generation, retrieve-or-reason decisions, review, document analysis, or
structured expansion inside an adaptive control loop, but their training
regimes differ substantially. Answer accuracy alone cannot show whether a
controller improved chain observability or merely spent more reasoning and
retrieval budget. At minimum, comparisons should match or expose retrieval
calls, tokens, stopping behavior, and complete-chain recall.

\subsection{Budgeted Evidence Selection and Context Construction}
\label{subsec:methods_utility}

Once a chain is present in $L_N$, selection determines whether it survives in
$E_K$. The key distinction is the decision unit: an independently scored
document, a list, a path, a set conditioned on partial evidence, or a
compressed representation. Table~\ref{tab:selection_design} summarizes the
consequences.

\begin{table}[t]
\caption{Selection mechanisms under a fixed reader budget.}
\label{tab:selection_design}
\centering
\scriptsize
\begin{tabularx}{\linewidth}{@{}p{2.2cm}p{2.0cm}p{2.6cm}p{2.7cm}Y@{}}
\toprule
Mechanism & Decision unit & Conditioning signal & What it can preserve & Main confound \\
\midrule
Pointwise reranking & one document & query--document relevance & individually relevant supports & low-surface bridge evidence \\
Listwise reranking & ordered candidate list & query plus candidate interactions & relative membership and order & membership and exposure changed together \\
Bridge / path selection & next node, passage, or path & partial chain or bridge entity & complementary chain elements & dependence on noisy intermediate state \\
Plan / subquestion selection & evidence per subtask & explicit decomposition or plan & step-aligned evidence & decomposition errors and duplicated evidence \\
Compression / notes & tokens, spans, or generated notes & salience or reader contribution proxy & more evidence within token budget & lost relation, provenance, or factual detail \\
\bottomrule
\end{tabularx}
\end{table}

Cross-encoders and neural rerankers, from passage reranking and
\textsc{monoT5} to \textsc{RankT5}, \textsc{RankLLaMA},
\textsc{RankGPT}, and \textsc{RankZephyr}~\cite{nogueira2019passage,nogueira2020monot5,zhuang2023rankt5,ma2024rankllama,sun2023rankgpt,pradeep2023rankzephyr},
primarily estimate query-conditioned relevance. Listwise language-model
rerankers~\cite{tang2024found} can model candidate interactions, but a
better listwise score is not yet evidence that a complementary bridge was
preserved. The required diagnostic is complete-chain preservation under the
same pool and final token budget.

Bridge- and path-aware methods explicitly condition on partial state.
\textsc{Baleen}~\cite{khattab2021baleen},
\textsc{HopRetriever}~\cite{li2020hopretriever},
\textsc{PathRetriever}~\cite{asai2020learning},
\textsc{BeamDR}~\cite{zhao2021beamdr},
\textsc{MuGER2}~\cite{yang2022muger}, and
\textsc{Bridge-R}~\cite{xu2024bridger} score next-hop candidates, bridge
entities, or chain hypotheses using more than the original query. Their design
is closer to the conditional contribution $u(e\mid q,C)$, but the score is
reader- and state-dependent and should not be presented as a universal notion
of relevance.

Decomposition and planning methods such as \textsc{DecompRC},
\textsc{Self-Ask}, \textsc{Least-to-Most}, \textsc{BREAK},
\textsc{Plan-and-Solve}, and
\textsc{ProgramFC}~\cite{min2019decomprc,press2022selfask,zhou2023leasttomost,wolfson2020break,wang2023planandsolve,pan2023programfc}
convert a complex information need into subqueries or executable steps. This
can make evidence selection more nearly single-hop, but it also introduces a
new failure source: a wrong or overly coarse decomposition can make the valid
chain unreachable before ranking begins.

Compression methods including \textsc{LLMLingua} and
\textsc{LongLLMLingua-2}~\cite{jiang2023llmlingua,pan2024longllmlingua2},
\textsc{RECOMP}~\cite{xu2023recomp}, \textsc{Chain-of-Note}~\cite{yu2023chainofnote},
and \textsc{xRAG}~\cite{cheng2024xrag} alter both the budget and the evidence
representation. A compressed context may retain an answer-bearing phrase
while removing the bridge relation or source provenance. Compression should
therefore report chain preservation and citation validity at matched tokens,
not only compression ratio and answer accuracy.

Supervision for selection can come from gold supports, path annotations,
reader-derived leave-one-out labels, verifier scores, or synthetic
counterfactual chains. Each signal has bias: gold supports are incomplete,
reader-derived contribution is model-specific, and verifier scores may reward
plausibility.
A strong study states the surrogate and reports whether its gain survives a
reader change.

\subsection{Evidence Presentation and Exposure}
\label{subsec:methods_exposure}

Exposure concerns how a fixed evidence set is made available to a particular
reader. Ordering and packing heuristics respond to documented position
sensitivity~\cite{liu2024lostmiddle,bai2024longbench,yen2024lcontext}; learned
approaches such as \textsc{PRCA}~\cite{yang2023prca} and position-robust
training~\cite{he2024positionrobust} alter the layout or reader. Fusion-in-
decoder architectures such as \textsc{FiD}, \textsc{FiD-Light},
\textsc{KG-FiD}, and \textsc{Atlas}~\cite{izacard2021fid,hofstatter2023fidlight,yu2022kgfid,izacard2022atlas}
reduce interference caused by concatenating all passages into one encoder
sequence, but they do not remove selection failure. Compression can also
improve access by shortening distractors, while simultaneously changing
membership at the token level.

The decisive experiment is a fixed-membership treatment: hold $E_K$, the
reader, and the total budget constant, then vary only ordering, grouping, or
packing. Without that control, an ``ordering'' gain may be a hidden selection
gain. Exposure is reader-specific, so the same layout should be tested across
at least the reader families for which a general claim is made.

\subsection{Cross-Evidence Fusion}
\label{subsec:methods_fusion}

Fusion methods act after usable evidence has reached the reader. Retrieve-
then-read systems such as \textsc{DrQA}, \textsc{ORQA}, \textsc{RAG},
\textsc{REALM}, \textsc{Atlas}, and \textsc{FiD}~\cite{chen2017drqa,lee2019orqa,lewis2020rag,guu2020realm,izacard2022atlas,izacard2021fid}
differ in retrieval training and passage encoding, but their fusion claim is
best isolated with oracle or controlled evidence. Graph-aware readers such as
\textsc{QA-GNN}, \textsc{GreaseLM}, \textsc{UniKGQA},
\textsc{ReaRev}, and \textsc{NSM}~\cite{yasunaga2021qagnn,zhang2022greaselm,jiang2023unikgqa,mavromatis2022rearev,he2021nsm}
constrain composition through explicit structure.

Decomposition and programmatic reasoning reduce the complexity of the reader
operation. \textsc{DecompRC}, \textsc{Self-Ask},
\textsc{Least-to-Most}, \textsc{Decomposed Prompting},
\textsc{Successive Prompting}, and \textsc{ReasonPath}~\cite{min2019decomprc,press2022selfask,zhou2023leasttomost,khot2023decomposed,dua2022successive,li2023reasonpath}
externalize intermediate steps; \textsc{PoT}, \textsc{PAL},
\textsc{Faithful CoT}, and code-based reasoning~\cite{chen2023pot,gao2023pal,lyu2023faithfulcot,wang2024executable}
make some operations executable. Sampling and search over reasoning traces,
including self-consistency, Tree-of-Thought, and Graph-of-Thought~\cite{wang2023selfconsistency,yao2023tot,besta2024got},
can increase the chance that one trajectory composes the evidence correctly,
but also add inference cost.

A fusion claim requires a chain-containing context and a declared
presentation policy. If accuracy remains low under oracle evidence, the
reader or decomposition is the binding failure. If accuracy improves only
when the retrieval pool changes, the result should not be attributed to
fusion alone.

\subsection{Verification and Causal Faithfulness}
\label{subsec:methods_faithfulness}

Verification methods differ in where their evidence comes from. Self-
verification and critique methods such as \textsc{Self-Refine},
\textsc{Reflexion}, \textsc{CoVe}, \textsc{SelfCheckGPT}, and self-
verification prompting~\cite{madaan2023selfrefine,shinn2023reflexion,dhuliawala2024cove,manakul2023selfcheckgpt,weng2023selfverification}
reuse the model's own judgments. Externally grounded approaches including
\textsc{CRITIC}, \textsc{RARR}, \textsc{Verify-and-Edit},
\textsc{CRAG}, \textsc{Self-Checker}, and attribution-oriented traces~\cite{gou2024critic,gao2023rarr,zhao2023verifyedit,yan2024crag,li2024selfchecker,li2023attribution}
retrieve or check evidence before revision. Process-supervised systems such
as the process reward model of Lightman et al., \textsc{Math-Shepherd},
\textsc{ReARTeR}, \textsc{RAG-Gym}, \textsc{HiPRAG}, and
\textsc{VERITAS}~\cite{lightman2024process,wang2023mathshepherd,sun2025rearter,xiong2025raggym,wu2025hiprag,xu2025veritas}
move step-level reasoning, search-decision, or faithfulness signals into
training. Their supervision targets differ: \textsc{HiPRAG} penalizes
suboptimal search decisions, whereas \textsc{VERITAS} rewards trace-level
faithfulness among information, reasoning, answers, and search actions.

Citation-generating systems such as \textsc{WebGPT}, \textsc{GopherCite},
and \textsc{ALCE}~\cite{nakano2021webgpt,menick2022gophercite,gao2023alce}
make attribution inspectable. Deletion tests~\cite{liu2023evaluating,xu2024docnli}
and bridge counterfactuals such as \textsc{CofCA}~\cite{wu2024cofca}
intervene on evidence more directly. These mechanisms should not be collapsed
into one ``faithfulness'' score: self-critique tests consistency, citation
metrics test entailment and coverage, and counterfactuals provide stronger
but still imperfect evidence of causal use.

\subsection{Hybrid and Agentic Systems: The Attribution Problem}
\label{subsec:methods_joint}

Hybrid systems combine search, selection, reasoning, and verification, so an
architectural name is least informative precisely where the system is most
capable. Table~\ref{tab:recent_methods_diagnostic} illustrates how recent
systems can be described without treating the coded stage as a performance
claim.

\begin{table}[t]
\caption{Representative recent systems and the diagnostics needed to support
their coded stage-aligned interpretation. The table describes mechanisms, not a
comparative performance ranking.}
\label{tab:recent_methods_diagnostic}
\centering
\scriptsize
\begin{tabularx}{\linewidth}{@{}p{2.1cm}p{2.7cm}p{2.2cm}Y@{}}
\toprule
System & Central control signal & Primary target (secondary effect) & Required supporting diagnostic \\
\midrule
Search-R1~\cite{jin2025searchr1} & RL-trained query generation during reasoning & observability & complete-chain recall and cost at matched calls \\
DeepRAG~\cite{guan2025deeprag} & retrieve-or-reason policy & observability (cost control) & retrieval frequency, chain coverage, and accuracy--cost curve \\
RAG-Gym~\cite{xiong2025raggym} & process supervision for search actions & observability (selection) & step-level rewards plus pool, context, and cost diagnostics \\
PAR-RAG~\cite{zhang2025parrag} & top-down plan--act--review loop & observability (fusion) & plan validity, chain recall, and review ablation \\
VERITAS~\cite{xu2025veritas} & fine-grained trace-faithfulness rewards & causal faithfulness & information--think, think--answer, and think--search faithfulness \\
MultiCube-RAG~\cite{shi2026multicuberag} & ontology-guided structured retrieval & observability & chain coverage, noise, and ontology-construction cost \\
\bottomrule
\end{tabularx}
\end{table}

The general rule is simple: when an end-to-end system improves answer
accuracy, the paper should show which observable object changed---the pool,
the final context, its presentation, the reader outcome under controlled
evidence, or the answer's response to evidence interventions. Otherwise a
gain from deeper search, more computation, or parametric answering is easily
misattributed to the advertised component.

\subsection{Cross-Family Synthesis}
\label{subsec:cross_family_synthesis}

Four cross-family findings follow from the diagnostic view. First, most
architectures that advertise multi-hop capability still make their strongest
claim at the acquisition boundary: iterative retrievers, graph expanders,
hierarchical indexes, and search agents all try to make an otherwise latent
chain reachable. Their results are only comparable when retrieval depth,
call count, and pool size are reported together.

Second, budgeted selection remains the least cleanly isolated middle stage.
Rerankers, decomposers, compressors, and graph pruners often change both
membership and ordering, so answer gains cannot by themselves show that a
valid chain survived the final budget. This is where complete-chain
preservation under matched passage or token budgets is most valuable.

Third, long context and graph structure help only when exposure is measured
separately from membership. A larger context, community summary, or graph
path can increase the chance that evidence is present while still placing
dependent units too far apart, overcompressing the bridge, or obscuring
provenance.

Fourth, causal faithfulness remains under-instrumented relative to answer
accuracy and citation quality. Verification and attribution methods make
support more inspectable, but deletion, conflict, and bridge-counterfactual
probes are needed before a paper can claim that the supplied chain caused the
answer rather than merely agreeing with it.

\section{Architectural Crosswalk}
\label{sec:architecture_view}

The diagnostic organization does not replace architecture. It clarifies why
architectural labels alone are insufficient for comparison. Table
\ref{tab:architecture_estimand} maps common families to the stages they most
often affect. The mapping is non-exclusive: an iterative retriever can also
perform selection, a graph system can support both acquisition and fusion,
and an agent can act across the full pipeline.

\begin{table}[H]
\caption{Architectural families and their typical evidence-chain effects. A
filled circle marks a common primary effect and an open circle a frequent
secondary effect.}
\label{tab:architecture_estimand}
\centering
\small
\begin{tabularx}{\linewidth}{@{}p{2.7cm}CCCCC@{}}
\toprule
Family & Observability & Selection & Exposure & Fusion & Faithfulness \\
\midrule
Single-shot retrieval & $\bullet$ & $\circ$ &  &  &  \\
Iterative / adaptive retrieval & $\bullet$ & $\bullet$ & $\circ$ &  &  \\
Graph-guided retrieval & $\bullet$ & $\bullet$ &  & $\circ$ & $\circ$ \\
Decomposition / planning & $\circ$ & $\bullet$ &  & $\bullet$ & $\circ$ \\
Compression / context packing &  & $\circ$ & $\bullet$ & $\circ$ &  \\
Fusion readers / structured reasoning &  &  & $\circ$ & $\bullet$ &  \\
Verification / attribution &  &  &  & $\circ$ & $\bullet$ \\
Agentic / RL search & $\bullet$ & $\bullet$ & $\circ$ & $\circ$ & $\circ$ \\
\bottomrule
\end{tabularx}
\end{table}

Three distinctions are especially important. First, graph-based methods are
not a single stage: graph construction can increase candidate coverage,
subgraph pruning can change selection, and message passing can change fusion.
Second, ``long context'' is not a substitute for retrieval analysis. It
changes the feasible budget and evidence granularity but can preserve or
worsen position sensitivity. Third, agentic systems add adaptive control
rather than a new diagnostic target. Their main evaluation problem is
attribution: additional calls can improve answer accuracy through deeper
search, better selection, more computation, or self-correction. Per-stage
and cost-normalized results are therefore more informative than a single
agent-versus-pipeline comparison.
\section{Benchmarks as Diagnostic Instruments}
\label{sec:benchmarks}

A benchmark is diagnostically useful when its annotations or interventions
make a failure stage observable while controlling preceding stages. Domain
and release year matter, but they are secondary to what the benchmark can
identify.

\subsection{Supporting Facts, Paths, and Chain Coverage}
\label{subsec:coverage_diagnostics}

\textsc{HotpotQA}, \textsc{2WikiMultiHopQA}, \textsc{IIRC},
\textsc{ComplexWebQuestions}, and \textsc{QAngaroo}~\cite{yang2018hotpotqa,ho2020wikimultihopqa,ferguson2020iirc,talmor2018complexwebquestions,welbl2018qangaroo}
provide varying degrees of support or path information. These labels allow
pool-level and context-level coverage diagnostics, but they identify
annotated chains rather than every valid chain or the chain actually used by
the model. \textsc{2WikiMultiHopQA}, \textsc{EntailmentBank},
\textsc{eQASC}, and \textsc{ProofWriter}~\cite{ho2020wikimultihopqa,dalvi2021entailmentbank,jhamtani2020eqasc,tafjord2021proofwriter}
make more intermediate structure explicit and are therefore better suited to
path or step analysis. Even then, a gold path need not be unique.

\subsection{Shortcut Resistance and Longer Dependencies}
\label{subsec:shortcut_diagnostics}

\textsc{MuSiQue}~\cite{trivedi2022musique} was designed to reduce
single-hop shortcuts in composed questions. \textsc{MoreHopQA}~\cite{schnitzler2024morehopqa}
stresses longer chains, while \textsc{HotpotQA-NoShortcut}~\cite{trivedi2020nohotpot}
and related filtered settings reduce examples that can be solved from one
support. \textsc{NovelHopQA}~\cite{novelhopqa2025} combines dependency with
longer narrative context. These datasets are useful for acquisition claims,
but a low end-to-end score still cannot reveal whether the chain was missing,
removed by the context budget, or mishandled by the reader.

\subsection{Heterogeneous Evidence and Long Context}
\label{subsec:hetero_diagnostics}

\textsc{HybridQA}, \textsc{OTT-QA}, \textsc{MultiModalQA},
\textsc{WebQA}, \textsc{ManyModalQA}, \textsc{TAT-QA},
\textsc{FeTaQA}, \textsc{VisDoM}, \textsc{DocHop-QA}, and
\textsc{WikiMixQA}~\cite{chen2020hybridqa,chen2021opentableqa,talmor2021multimodalqa,chang2022webqa,hannan2020manymodalqa,zhu2021tatqa,nan2022fetaqa,wang2024visdom,suri2024dochop,he2024wikimixqa}
require chains that cross text, tables, documents, or images. Their primary
value is to test whether evidence representation and fusion remain coherent
across substrates; alignment errors can dominate ordinary retrieval error.

Long-context suites such as \textsc{LongBench}, $\infty$Bench,
\textsc{HELMET}, \textsc{Loong}, and \textsc{ZeroSCROLLS}~\cite{bai2024longbench,zhang2024infinitybench,yen2024helmet,wang2024loong,shaham2023zeroscrolls}
can reveal access failures, but their default protocols rarely isolate
presentation from membership. To support an exposure claim, a study should
construct fixed-set variants that move the same support chain across
positions or separations while holding tokens and reader constant.

\subsection{Support, Attribution, and Counterfactual Use}
\label{subsec:faithfulness_diagnostics}

Fact-verification datasets such as \textsc{FEVER}, \textsc{HoVer},
\textsc{EX-FEVER}, \textsc{SciFact}, and \textsc{AVeriTeC}~\cite{thorne2018fever,jiang2020hover,ma2024exfever,wadden2020scifact,schlichtkrull2024averitec}
measure whether claims are supported or refuted by evidence. Citation-oriented
benchmarks such as \textsc{ALCE}~\cite{gao2023alce}, together with
\textsc{HaluEval} and \textsc{ExpertQA}~\cite{li2023halueval,malaviya2024expertqa},
measure attribution, factuality, or support quality. These are important
faithfulness proxies, but support entailment does not show that the evidence
caused the answer. Bridge-counterfactual protocols such as
\textsc{CofCA}~\cite{wu2024cofca}, conflict tests, and carefully controlled
deletion experiments provide stronger evidence because they intervene on the
chain.

\subsection{End-to-End RAG and Agentic Behavior}
\label{subsec:rag_specific_diagnostics}

\textsc{MultiHop-RAG}, \textsc{FRAMES}, \textsc{CRAG},
\textsc{FanOutQA}, and \textsc{RAGTruth}~\cite{tang2024multihoprag,krishna2024frames,yang2024crag,zhu2024fanoutqa,niu2024ragtruth}
measure end-to-end RAG behavior under complex information needs.
\textsc{ToolBench}, \textsc{AgentBench}, and \textsc{WebArena}~\cite{qin2023toolbench,liu2024agentbench,zhou2024webarena}
add tool or environment interaction. They are valuable deployment tests, but
their scores combine action policy, retrieval affordances, context
construction, reader reasoning, and cost. Stage-level logs are required for
mechanism claims.

\begin{table}[t]
\caption{Benchmark families interpreted as diagnostic instruments.}
\label{tab:benchmark_families_extended}
\centering
\scriptsize
\begin{tabularx}{\linewidth}{@{}p{2.5cm}p{2.8cm}p{2.3cm}Y@{}}
\toprule
Family & Available signal & Best-identified target & Principal blind spot \\
\midrule
Supporting-fact QA & document or sentence supports & observability and selection & incomplete chains and shortcut answers \\
Path / proof annotated & ordered paths, trees, or steps & fusion and process validity & alternative valid derivations \\
Shortcut-resistant QA & composed or filtered questions & dependent acquisition & stage of failure remains mixed \\
Heterogeneous evidence & table, text, graph, or image links & cross-substrate fusion & alignment and representation confounds \\
Long-context suites & long inputs and position variation & exposure when membership is controlled & default tasks mix selection and layout \\
Verification / citation & support labels, entailment, citations & support validity and attribution & causal use is not identified \\
Counterfactual / conflict & altered bridge or evidence facts & causal faithfulness & intervention naturalness and redundancy \\
Agentic / tool use & action traces and environment outcomes & joint control and efficiency & tool affordances dominate comparisons \\
\bottomrule
\end{tabularx}
\end{table}

\subsection{Metric--Claim Alignment}
\label{subsec:metric_alignment}

Answer EM and F1 are deployment-level metrics rather than stage-aligned
metrics.
Supporting-fact F1 approximates annotated membership; Top-$N$ recall estimates
pool observability; citation precision estimates attribution; and
counterfactual sensitivity estimates response to an intervention. None is a
universal ``multi-hop RAG quality'' score.

\begin{table}[t]
\caption{Metrics and the claims they can support when the required control is
satisfied.}
\label{tab:metric_estimand_alignment}
\centering
\scriptsize
\begin{tabularx}{\linewidth}{@{}p{2.5cm}p{3.0cm}p{2.8cm}Y@{}}
\toprule
Claim & Suitable metric & Conditional subset / control & Misleading substitute \\
\midrule
Pool contains a valid chain & complete support or path recall at Top-$N$ & fixed corpus, evidence unit, and retrieval budget & answer F1 alone \\
Budget preserves the chain & selected-chain recall at $K$ & chain present in pool; fixed token budget & reranker relevance score \\
Presentation improves access & fixed-set ordering or packing gap & identical membership and reader & reranking that also changes membership \\
Reader composes evidence & oracle- or controlled-evidence EM/F1; step validity & chain supplied and presentation declared & Top-$N$ retrieval recall \\
Answer is faithfully supported & joint citation, deletion, conflict, and counterfactual diagnostics & answer-correct and chain-present subsets & citation overlap alone \\
System is efficient & tokens, calls, latency, cost--quality curve & matched model/API/hardware conditions & selector-only latency \\
\bottomrule
\end{tabularx}
\end{table}

A balanced empirical portfolio should include at least one benchmark with
chain annotations, one that reduces shallow shortcuts or increases dependency
length, and one that tests support or counterfactual behavior. No individual
dataset substitutes for conditional evaluation on the system's own pool,
context, ordering, reader, and evidence interventions.

\section{Faithfulness as Causal Evidence Use}
\label{sec:causal_faithfulness}

Correctness, attribution, and causal use are distinct. A model can answer
correctly without citing evidence, cite an entailing passage that it did not
use, or produce a well-supported chain whose final composition is wrong.
Faithfulness evaluation should therefore ask how the answer responds when the
relevant evidence is withheld, contradicted, or altered, rather than treating
a plausible citation as proof of use.

\subsection{Operational Probes}
\label{subsec:faithfulness_probes}

For a candidate chain $C\subseteq E_K$, a simple deletion statistic is
\begin{equation}
D_{\mathrm{del}}(C)=
P(a^\star\mid q,\pi(E_K))-P(a^\star\mid q,\pi(E_K\setminus C)).
\end{equation}
A large value suggests that $C$ is load-bearing, but it is not a definition of
faithfulness. Deletion can create an unnatural input, and redundant valid
chains can make a faithful answer insensitive to any one removal.

\begin{table}[t]
\caption{Complementary probes of evidence use. No single row identifies
causal faithfulness in all settings.}
\label{tab:faithfulness_probes}
\centering
\scriptsize
\begin{tabularx}{\linewidth}{@{}p{2.4cm}p{2.8cm}p{2.7cm}Y@{}}
\toprule
Probe & What it asks & Strong evidence when & Failure mode \\
\midrule
Closed-book / conflict control & does the answer follow supplied evidence rather than parametric memory? & evidence conflicts are coherent and answerable & model may distrust artificial conflicts \\
Deletion / sufficiency & does removing a support chain reduce the answer probability? & input remains natural and alternatives are controlled & redundancy and distribution shift \\
Bridge counterfactual & does an altered intermediate fact propagate to the answer? & replacement yields a coherent alternative chain & expensive construction and leakage from unchanged cues \\
Citation entailment / coverage & do cited units support the generated claims? & claim segmentation and entailment judgments are reliable & post-hoc citation without causal use \\
\bottomrule
\end{tabularx}
\end{table}

Counterfactual bridge substitution is often stronger than deletion because it
preserves a coherent context while changing the fact that should propagate
through the chain. Closed-book and evidence-conflict controls estimate the
availability and dominance of parametric knowledge. Citation entailment is
necessary for auditable output, but it measures support compatibility rather
than the internal cause of generation. These probes should be combined rather
than averaged into an opaque single score.

\subsection{Reporting Faithfulness}

A credible audit reports at least three joint quantities on the answer-correct
subset: citation entailment or support precision, sensitivity to a declared
evidence intervention, and closed-book or conflict behavior. Results should
also be stratified by whether an annotated chain is present in the final
context. Faithfulness on chain-absent examples mainly measures parametric
answering or retrieval failure; faithfulness averaged over all examples can
hide that distinction.

Verification methods should be described by the stage at which they act.
Self-critique tests internal consistency, externally grounded verification
checks or retrieves additional support, process supervision changes the
training objective, and counterfactual training changes the data distribution.
Their outputs are complementary evidence about causal use, not interchangeable
measurements of one universal property.

\section{Open Problems and Tradeoffs}
\label{sec:tradeoffs}

\paragraph{Dependent retrieval under controlled cost.}
Iterative and agentic search can expose bridge evidence that is invisible to
the original query, but each step also increases latency, distractor mass, and
query-drift risk. Progress requires accuracy--cost curves, chain-recall curves
by call count, and stopping policies trained against evidence acquisition
rather than only final-answer reward.

\paragraph{Set- and path-level selection.}
Most rerankers still score documents independently or listwise against the
original query. Multi-hop selection requires complementarity: a passage can
be weak alone and decisive with another passage. Core-conditioned scoring,
path-level objectives, set selection, and counterfactual contribution labels
are promising, but they must be evaluated under fixed pools and reader token
budgets.

\paragraph{Evidence granularity, compression, and provenance.}
Passages, sentences, graph triples, summaries, and generated notes induce
different chain lengths and costs. Compression can preserve an answer string
while removing the relation or provenance required to justify it. Studies
should report the evidence unit, token accounting, chain preservation, and
citation recoverability after compression.

\paragraph{Reader-specific context organization.}
Long-context models do not make presentation irrelevant. Ordering, separation,
grouping, and instruction layout interact with the reader architecture and
training distribution. Exposure metrics should therefore be reader-specific,
and presentation claims require fixed-membership tests across the reader
families to which the claim is intended to generalize.

\paragraph{Identifiable benchmarks with alternative chains.}
Many benchmarks annotate one support path even when several valid chains
exist. New evaluation instruments should include alternative chain labels,
controlled distractor density, fixed-set ordering variants, oracle contexts,
and coherent bridge counterfactuals. These additions are more informative
than another end-to-end leaderboard with no failure localization.

\paragraph{Attribution and robustness in adaptive systems.}
Agents, graph builders, and verifiers increase both capability and attack
surface. Retrieval-side prompt injection, poisoned or stale corpora, source
conflicts, and unverifiable generated graph edges can corrupt the chain while
leaving the answer fluent. Trust-aware retrieval, provenance-preserving
context construction, and per-action audit logs should be evaluated alongside
accuracy and efficiency.

\section{Reporting Guidelines for Multi-Hop RAG}
\label{sec:reporting_guidelines}

The evidence-chain framework implies a compact reporting protocol. A paper
need not optimize every stage, but it should measure the stage it claims to
improve, expose major spillovers, and state the budget under which the claim
holds.

\begin{table}[t]
\caption{Minimum evidence-chain reporting for multi-hop RAG studies.}
\label{tab:minimum_reporting}
\centering
\footnotesize
\begin{tabularx}{\linewidth}{@{}p{2.8cm}p{2.1cm}Y@{}}
\toprule
Reported quantity & Claim & Required control \\
\midrule
Chain/path recall at Top-$N$ & pool observability & fixed corpus, retrieval budget, and evidence unit \\
Chain preservation at budget $K$ & selection & chain present in the pool; fixed token budget \\
Order/packing gap on a fixed set & exposure & identical evidence membership, budget, and reader \\
Controlled-context EM/F1 & fusion & valid chain supplied and presentation specified \\
Entailment plus intervention probes & causal faithfulness & answer-correct and chain-present subsets \\
Tokens, calls, and end-to-end latency & efficiency & comparable model, API, and hardware conditions \\
\bottomrule
\end{tabularx}
\end{table}

Four practices are necessary for interpretable comparison. First,
retrieval depth $N$, final context budget $K$, evidence granularity, reader
model, and maximum calls must appear next to the headline score. Second,
stage-aligned metrics must be reported on the corresponding conditional subset:
selection on chain-present pools, exposure on fixed membership, fusion on
controlled evidence, and faithfulness on answer-correct cases with available
support. Third, hybrid and agentic systems need stage-aligned ablations and
cost--quality curves because additional computation can conceal the source of
an improvement. Fourth, benchmark and survey papers should publish their
coding assumptions and distinguish verified records from discovery records.

Incomplete gold chains make recall a lower-bound proxy; reader-derived
contribution is model-specific; deletion sensitivity is affected by
redundancy; and citation entailment does not establish causal use. Stating
these limitations produces more comparable evidence than assigning one
overloaded score to ``RAG quality.''

\section{Conclusion}
\label{sec:conclusion}

Multi-hop RAG is an information-retrieval problem over dependent evidence,
not simply a larger retrieve--then--read pipeline. The evidence-chain view
separates five questions that end-to-end answer accuracy conflates: whether a
valid chain enters the pool, survives the reader budget, is organized for
access, is composed correctly, and faithfully supports the answer. This
organization connects sparse and dense retrieval, graph expansion,
decomposition, compression, fusion readers, verification, and agentic search
without treating an architectural family as an explanation of performance.

The main methodological implication is conditional evaluation. Pool recall
cannot support a reader claim; reranking that changes both membership and
order cannot isolate exposure; oracle-evidence accuracy does not establish
faithfulness; and citation overlap does not prove causal use. Each claim
requires a corresponding observable object, subset, intervention, and budget.
The companion artifact operationalizes the same discipline for the survey
itself by separating discovery breadth from source-verified taxonomy records.

Future progress will depend less on adding another pipeline label and more on
building controlled instruments for dependent acquisition, complementary
selection, reader-specific context organization, cross-evidence composition,
and counterfactual support. These instruments are what allow improvements to
be compared across rapidly changing retrievers, readers, graph indexes, and
search agents.

\section*{Data and Artifact Availability}

The companion artifact is available at
\url{https://github.com/TsingyuL/multi-hop-rag-survey}. It contains the search
protocol, machine-readable method and benchmark records, coding fields,
source URLs, status labels, figure sources, and validation scripts. Versioned
releases allow the manuscript, bibliography, figures, and coded records to be
associated with the same frozen snapshot.

\appendix

\section{Artifact Status and Update Policy}
\label{app:coverage_counts}

The artifact separates three layers that serve different purposes.
\emph{Discovery records} maximize recall during literature search and may
contain adjacent or unverified work. \emph{Source-verified records} have been
checked against the canonical paper or official project page and can support
artifact-derived taxonomy analyses. \emph{Manuscript citations} additionally
include background and component papers needed for exposition. These layers
overlap and should not be interpreted as a PRISMA-style screening funnel.

Each public release records its literature cut-off, canonical-source
policy, status vocabulary, and correction history. Numerical catalog summaries
are generated from the release rather than copied into the source by hand,
which prevents stale counts and restricts artifact-derived tables to records
that satisfy the stated verification rule.

\section{Screening and Coding Checklist}
\label{app:protocol}

A method is promoted to reviewed status only after the following fields have
been checked against the full paper or an official project page: canonical
title and venue/version; multi-hop scope; evidence substrate; retrieval and
reader budgets; architectural family; primary and secondary evidence-chain
stages; benchmark and metric claims; source and code URLs; and a short passage
supporting the taxonomy assignment. Comparative tables should be generated
from reviewed records rather than maintained manually.

\input{main_TOIS_verified.bbl}

\end{document}
